\chapter{Appendix A: Muscle Force and Joint Torque Calculations}\label{app:muscle_calcs}
\graphicspath{{figs/92-AppendixA/}}

This appendix documents the reference implementation of the force, torque, and compliance equations used in Chapters~\ref{ch:methods}--\ref{ch:results}. When present in Matlab code, file names are given so that any number in this dissertation can be traced to the code that produced it.

    % ---------------------------------------------------------------------------
% muscle_model_appendix_draft.tex — appendix-ready draft covering the two
% muscle formulations used in this work (MuJoCo muscle actuator as produced
% by the conversion, and OpenSim's Thelen2003/Millard2012 Hill-type models),
% and the quantitative force comparison along the subject01 walking
% trajectory. Written 2026-09-14; regenerate numbers with
% Code\MuJoCo_SNS\spinal\_muscle_force_compare.py (muscle_force_compare.csv,
% figures\muscle_force_compare.png).
% ---------------------------------------------------------------------------
\section{Human muscle force calculations}
\subsection{OpenSim Thelen2003 and Millard2012 Formulations}

OpenSim's Thelen2003Muscle comprises a contractile element (activation
dynamics as above; active force--length curve; Hill-type force--velocity
curve) in series with a compliant tendon, plus a parallel passive fiber
element, with fiber pennation. Unlike the MuJoCo actuator, the fiber
length is \emph{not} the musculotendon length: at every instant the
model solves the force equilibrium
$F_{\text{tendon}}(L_{\text{mt}} - L_{\text{fiber}}\cos\alpha)
 = F_{\text{fiber}}(L_{\text{fiber}}, \dot{L}_{\text{fiber}}, a)$,
so tendon compliance filters and delays force transmission and stores
energy. The tendon is a quadratic-then-linear ``strip'' curve anchored
at the slack length $L_s$ (force zero until $L_{\text{mt}}$ stretches
the tendon beyond $L_s$). Millard2012EquilibriumModel, OpenSim's newer
default, retains the same equilibrium structure but replaces the curve
shapes (an asymmetric Gaussian active-force--length curve with a
plateau, an exponentially stiffening tendon) and adds fiber damping; it
is numerically better behaved at extreme lengths. Both OpenSim models
therefore differ from the converted MuJoCo model in the same single
respect: series-elastic compliance.

\section{Muscle Model Formulations and Cross-Platform Force Comparison}
\label{sec:muscle-models}

\subsection{MuJoCo Muscle Actuator (as converted)}

The conversion from OpenSim to MuJoCo re-expresses each Thelen2003
muscle as a MuJoCo \texttt{general} actuator with the built-in muscle
dynamics. Activation is a first-order process with distinct rise and
fall time constants (0.01\,s and 0.04\,s in the converted model). Force
is computed from the actuator's normalized length
$\bar{L} = (L - L_{\min}) / (L_{\max} - L_{\min})$, where the length
range is taken from the converted geometry, as
\begin{equation}
  F = F_{\max}\,\bigl[\,a\, f_L(\bar{L})\, f_V(\dot{L}) - f_P(\bar{L})\,\bigr],
  \label{eq:mj-muscle}
\end{equation}
with $a \in [0,1]$ the activation, $f_L$ a force--length curve that
peaks inside the normalized range, $f_V$ a force--velocity curve, and
$f_P$ a passive element. The peak force scale $F_{\max}$ (actuator
parameter \texttt{gainprm[2]}) is set by the converter such that the
peak force over the model's range of motion matches the OpenSim
\texttt{max\_isometric\_force}; audited across all 84 active muscles,
the deviation is $\le 15\%$ (Sec.~\ref{sec:muscle-models-compare}). Two
properties matter for interpretation: (i) the tendon is \emph{rigid} --
there is no series-elastic element and no tendon slack length; force at
the joint equals the contractile-element force immediately; and (ii)
all curve parameters are shared defaults of the converted model rather
than per-muscle fits.

\subsection{Cross-Platform Force Comparison}
\label{sec:muscle-models-compare}

Both models were driven with identical kinematics and identical
activations along the subject01 walking trajectory (the same IK frames
used for the static-optimization back-solve). The OpenSim side is the
StaticOptimization solution's own muscle forces (full Thelen dynamics
with compliant tendon); the MuJoCo side replays the identical joint
trajectory with those activations. Table~\ref{tab:force-r2} reports, per
muscle, the coefficient of determination $R^2$ between the two force
time series, the shape correlation, and the best-aligning lag.

\begin{table}[h]
  \centering
  \caption{MuJoCo vs.\ OpenSim muscle force along the subject01 walk at
  matched activations. Lag in analysis frames ($\approx 17$\,ms each).}
  \label{tab:force-r2}
  \small
  \begin{tabular}{lrrr}
    \toprule
    Muscle & $R^2$ & shape corr$^2$ & lag [frames] \\
    \midrule
    soleus   & 0.742 & 0.800 & 1 \\
    med.\ gastrocnemius & 0.888 & 0.894 & 2 \\
    lat.\ gastrocnemius & 0.884 & 0.887 & 2 \\
    tibialis anterior & 0.945 & 0.945 & 0 \\
    vastus lateralis & 0.941 & 0.945 & 1 \\
    rectus femoris & $-1.15$ & 0.384 & 0 \\
    semimembranosus & 0.565 & 0.621 & 0 \\
    biceps femoris short & $-0.26$ & 0.868 & 0 \\
    psoas & 0.927 & 0.934 & 1 \\
    gluteus maximus (2) & 0.926 & 0.943 & 1 \\
    tensor fasciae latae & 0.933 & 0.935 & 0 \\
    peroneus brevis & 0.977 & 0.989 & 0 \\
    \bottomrule
  \end{tabular}
\end{table}

For ten of twelve representative muscles the two formulations agree
with $R^2 = 0.87$--$0.98$ and mean-force ratios within $\pm 22\%$, with
\emph{no} systematic phase shift (best-aligning lag 0--2 frames,
$\le 34$\,ms). The absence of a phase lead in the rigid-tendon model
indicates that, along physiological trajectories, missing tendon slack
compliance does not desynchronize force transmission; its signature is
instead visible in waveform \emph{shape}: for the soleus (the largest
series-elastic energy store in human gait), the OpenSim compliant
tendon carries substantial force through mid-stance while the MuJoCo
rigid-tendon force follows the activation dips --- the expected
smoothing of the series elastic element. The two outliers are
anatomically explained: rectus femoris was rerouted in the conversion
over the vastii's patella-tracking via point, and its moment-arm
polyfit deviates from OpenSim's path (the same coupler artifact noted
for the vastii); biceps femoris short head agrees in shape
($r^2 = 0.87$) but carries a scale offset. We conclude that the
converted model reproduces OpenSim muscle force generation along
physiological trajectories to a degree sufficient for neural-control
studies, and that if tendon-elasticity phenomena (energy storage,
force smoothing) become material to a claim, the appropriate upgrade
is a compliant-tendon formulation rather than parameter retuning of
the rigid-tendon model.

% \begin{figure}[h]
%   \centering
%   \includegraphics[width=0.9\textwidth]{muscle_force_compare.png}
%   \caption{Normalized muscle force along the analyzed walk at matched
%   OpenSim static-optimization activations (MuJoCo: rigid tendon;
%   OpenSim: Thelen compliant tendon). Soleus illustrates the signature
%   of the missing series compliance: OpenSim force is carried through
%   mid-stance by tendon stretch while MuJoCo force follows the
%   activation.}
%   \label{fig:force-compare}
% \end{figure}

    


\section{BPA Force Model Implementation}

\subsection{Strain Measures}
The musculotendon path length $l_{LMT}$ is the sum of straight segments between route via points, transformed into a common frame at each frame crossing. The contractile strain used by the force model subtracts the artificial tendon and both end-cap fittings from the path length before normalization by the resting length,
\begin{equation}
    \varepsilon \;=\; \frac{l_{0} - \left(l_{LMT} - l_{ten} - 2\,l_{fitting} - \chi_{0}\right)}{l_{0}},
    \label{eq:app_strain}
\end{equation}
where $\chi_{0}$ is the constant length offset of Section~\ref{sec:improved_model} (zero in the uncorrected model). [Provenance: \texttt{MonoPamDataExplicit.m} get.Contraction, lines 210--217; \texttt{minimizeFlxPin.m} lines 150--164. End-cap fitting length $l_{fitting}=\qty{25.4}{\mm}$ is a constructor default, \texttt{MonoPamDataExplicit.m} line 69.]

The maximum free-load contraction at \qty{620}{\kPa} is $\varepsilon_{620} = (l_{0}-l_{620})/l_{0}$, where $l_{620}$ is the BPA length at \qty{620}{\kPa} with no external load (code variable \texttt{KMAX}; \texttt{MonoPamDataExplicit.m} line 278). The relative strain is $\varepsilon^{*} = \varepsilon / \varepsilon_{620}$ (line 281).

\subsection{Normalized Force Surface}
The normalized isometric force surface of Equation~(\ref{eq:Fstar}) is implemented in closed form in \texttt{Functions/f\_festo.m} (lines 1--26),
\begin{equation}
    F^{*} = c_{0}\left(e^{-c_{1}\varepsilon^{*}} - 1\right)
            + P^{*}\,e^{-c_{2}(\varepsilon^{*})^{2}},
    \qquad F^{*}\le 0 \;\Rightarrow\; 0,
    \label{eq:app_fstar}
\end{equation}
with $P^{*} = P/\qty{620}{\kPa}$. The zero-clamp $F^{*}\le 0 \Rightarrow 0$ is applied by the callers (e.g.\ \texttt{Functions/festo4.m} line 35), not inside \texttt{f\_festo.m} itself. The coefficients were fitted with \texttt{fittype} nonlinear least squares and least-absolute-residual robustness (\texttt{Testing\_Data/bpaFits.m}, lines 45--60 and 115--140, lower bounds $[0,0,0]$) and are reported in Table~\ref{tab:Fstar}. The code additionally holds a 40\,mm fit not used in this dissertation ($c_0 = 0.1224$, $c_1 = 10.47$, $c_2 = 2.023$; \texttt{f\_festo.m} lines 17--20).

At runtime the torque model may alternatively evaluate a spline interpolation of the manufacturer's datasheet curves (\texttt{Functions/festo4.m}, backed by \texttt{FestoLookup.mat}), with force clamped to zero for $\varepsilon^{*}>1$ or $F^{*}<0$. Equation~(\ref{eq:app_fstar}) replaced this lookup in the corrected model.

\subsection{Maximum Force}
The resting-length-dependent maximum force of Equation~(\ref{eq:Fmax_3D}) is implemented in \texttt{Functions/maxBPAforce.m} (10\,mm: line 46; 20\,mm: line 57; duplicate for 20\,mm in \texttt{Functions/Fmax20.m}) with the fixed \qty{7.5}{\mm} end-contact offset. The fitted 20\,mm coefficients carry confidence intervals $a_1 \in [1.4745, 1.5011]$\,\si{\N\per\kPa} and $a_2 \in [0.0238, 0.0258]$\,\si{\per\kPa\per\metre} (\texttt{Fmax20.m} lines 26--27). Total force is the product $F = F^{*}(\varepsilon^{*},P^{*})\,F_{620}(l_{0})$, resolved into a vector along the force direction (\texttt{MonoPamDataExplicit.m} lines 283--287).

\section{Muscle Path Length and Moment Arm}

Path length accumulates the Euclidean norm of each segment between via points, applying the body-frame transformation whenever the path crosses between the femur and tibia frames (\texttt{MonoPamDataExplicit.m} lines 78--95, total at lines 122--133).

The vector moment arm is implemented as the perpendicular foot of the joint origin onto the line of action,
\begin{equation}
    \vec{r} \;=\; \vec{p} \;-\; \hat{u}\left(\hat{u}\cdot\vec{p}\right),
    \label{eq:app_perpfoot}
\end{equation}
where $\vec{p}$ is the vector from the joint origin to the muscle-path crossing point and $\hat{u}$ is the unit force direction (\texttt{MonoPamDataExplicit.m} lines 156--168; identical implementations in \texttt{MonoMuscleData.m} lines 111--123, \texttt{minimizeFlxPin.m} lines 400--409, and \texttt{MonoPam\_mult.m} lines 186--201). This is the same quantity as the cross-product projection of Equation~(\ref{eq:project_ma}); the planar scalar moment arm about the joint axis is the magnitude of the in-plane components, computed as \texttt{hypot(mA(:,1), mA(:,2))} (\texttt{predictKneeFlexor20mm.m} line 199). Torque about the joint $z$-axis is $\vec{M} = \vec{r}\times\vec{F}$, evaluated only where the BPA is in tension (values are discarded for $\varepsilon < -0.02$, the excessive stretching region; \texttt{MonoPamDataExplicit.m} line 308).

\section{Compliance Model Implementation}

The compliance-corrected muscle length entering Equation~(\ref{eq:Lm}) is computed as $l_{m} = l_{LMT} - \chi_{0} - l_{ten} - 2\,l_{fitting}$ (\texttt{minimizeFlxPin.m} line 254). In the two-bracket fit of record, the insertion-bracket stiffness matrix and the projected scalar stiffness along the force direction are
\begin{equation}
    \mathbf{K}_{br} = \mathrm{diag}(\chi_{1},\,\chi_{2},\,\chi_{1}),
    \qquad
    k_{\hat{u},br} = \left(\hat{u}^{\,T}\mathbf{K}_{br}^{-1}\hat{u}\right)^{-1},
    \label{eq:app_kbr}
\end{equation}
matching Equation~(\ref{eq:bktstiffness}), and the origin-side second bracket contributes an additional scalar compliance $c_{b2}$ in series (Appendix~\ref{app:model_geometry}, Table~\ref{tab:app_stiffness_arrays}) [Provenance: \texttt{minimizeFlxPin2brk.m} lines 387--418]. The earlier single-bracket evaluator used $\mathrm{diag}(\chi_{1},\chi_{2},\chi_{2})$ and is superseded; its array survives only in the legacy file (\texttt{minimizeFlxPin.m} lines 266--272).

The artificial tendon is modeled as an axial spring with stiffness
\begin{equation}
    k_{ten} \;=\; \frac{n\,A_{c}\,E}{L_{ten}},
    \label{eq:app_tendon}
\end{equation}
with $n=2$ parallel cable strands ($\times 6$ for the 20\,mm assembly), $A_{c} = \qty{1.51e-6}{\metre\squared}$ (19-strand cable), and $E = \qty{193e9}{\pascal}$ (\texttt{minimizeFlxPin.m} lines 460--471; \texttt{MonoPam\_mult.m} lines 606--617). Series combination with the projected bracket compliance gives $k_{eq} = \left(k_{\hat{u},br}^{-1} + k_{ten}^{-1}\right)^{-1}$ (lines 273--275).

The scalar compliance displacement $\delta$ of Equation~(\ref{eq:complianceforcebalance}) is found numerically with \texttt{fzero} on the bracket $[0,\, l_{m}-l_{620}]$,
\begin{equation}
    F_{620}(l_{0})\,F^{*}\!\left(\varepsilon^{*}(\delta),P^{*}\right) - k_{eq}\,\delta = 0,
    \label{eq:app_forcebalance}
\end{equation}
with $\delta \equiv 0$ when the uncorrected relative strain already exceeds unity (\texttt{minimizeFlxPin.m} lines 291--305). The vector bracket deflection is then $\boldsymbol{\epsilon}_{br} = \mathbf{K}_{br}^{-1}\left(\norm{\vec{F}}\hat{u}\right)$ (line 324), the cable stretch is $\norm{\vec{F}}/k_{ten}$ (line 334), and the attachment point, path, moment arm, and torque are recomputed. For the two-BPA flexor grouping, the coupled balance shares one common force between both muscles and is solved analogously on $[0, F_{620}]$ (\texttt{MonoPam\_mult.m} lines 490--564).

\section{Wrapping-Loss Implementation}

For the pinned-knee extensor, Equation~(\ref{eq:deltaL}) is implemented segment-wise over the wrap arc with radii $R_1 = \qty{12}{\mm}$ and $R_2 = \qty{40}{\mm}$ and segment boundaries $\theta_{1} = \ang{27}$ and $\theta_{2} = \ang{-23}$ (\texttt{minimizeExtX3.m}, nested \texttt{Contraction}, lines 327--370). For the biomimetic-knee 20\,mm cell, the arc-length factor is the routed bend measure $b(\theta_{k})$ (arc length swept by the wrap, built with the \texttt{wRap} bend radius in \texttt{buildKneeFlexorRoute20mm.m} lines 265--268), giving $\Delta l = \chi_{3}\, b(\theta_{k})\,(1-\varepsilon^{*})^{2}$ (\texttt{MonoPam\_mult.m} lines 398--441), with a hardcoded-geometry fallback for degenerate routes (lines 428--434). The loss enters only the force-producing strain; the kinematic (prediction) strain excludes it (\texttt{minimizeExtX3.m} lines 236--240).

